\documentclass[11pt]{article}
% ======================================================================
%  Preambule commun - Sujets Bac Serie A (Mathematiques)
%  Style sobre : noir et blanc, sans couleurs, sans filigrane,
%  sans en-tetes ni pieds de page decoratifs.
% ======================================================================
\usepackage[utf8]{inputenc}
\usepackage[T1]{fontenc}
\usepackage{lmodern}
\usepackage{textcomp}
\usepackage{amsmath,amssymb}
\usepackage{enumitem}
\usepackage{array}
\usepackage[a4paper,margin=2.3cm]{geometry}
\usepackage{fancyhdr}
\usepackage{tikz}

% Symbole degre robuste + justification souple
\DeclareUnicodeCharacter{00B0}{\ensuremath{^\circ}}
\tolerance=2000
\emergencystretch=2em
\frenchspacing

% En-tete / pied de page minimalistes (numero de page seul)
\pagestyle{fancy}
\fancyhf{}
\renewcommand{\headrulewidth}{0pt}
\fancyfoot[C]{\thepage}

% Macros mathematiques
\newcommand{\R}{\mathbb{R}}
\newcommand{\N}{\mathbb{N}}
\newcommand{\Z}{\mathbb{Z}}
\newcommand{\vi}{\vec{\imath}}
\newcommand{\vj}{\vec{\jmath}}

% Titres de blocs
\newcommand{\exo}[2]{\bigskip\noindent{\large\bfseries Exercice #1}\hfill\textit{#2}\par\nopagebreak\smallskip}
\newcommand{\partie}[1]{\medskip\noindent{\bfseries Partie #1}\par\nopagebreak\smallskip}
\newcommand{\entetesujet}[1]{\begin{center}{\Large\bfseries #1}\end{center}\vspace{0.3em}\hrule\vspace{1.1em}}

% Questions numerotees : niveau 1 = chiffres gras, niveau 2 = a. b. c.
\newlist{questions}{enumerate}{1}
\setlist[questions]{label=\textbf{\arabic*.},leftmargin=1.8em,itemsep=3pt,topsep=3pt}
\newlist{sousq}{enumerate}{1}
\setlist[sousq]{label=\textbf{\alph*.},leftmargin=1.8em,itemsep=2pt,topsep=2pt}


\begin{document}

\entetesujet{Sujet bac 2017 -- Série A}

\exo{1}{8 points}
Soit $P$ le polynôme défini sur $\R$ par :
\[ P(x) = 3x^3 - 8x^2 - x + 10 \]
\begin{questions}
  \item Calculer $P(-1)$ et $P(2)$.
  \item Développer l'expression $(x - 2)(x + 1)(3x - 5)$.
  \item En déduire les solutions de l'équation $P(x) = 0$.
  \item Résoudre dans $\R^2$ le système $\begin{cases} x + 2y = 0 \\ 2x - y = 5 \end{cases}$
  \item En déduire les solutions, dans $\R^2$, du système $\begin{cases} \ln x + 2\ln y = 0 \\ 2\ln x - \ln y = 5 \end{cases}$
  \item On donne le développement décimal du réel $b$ : $b = 1{,}123\,123\,123\cdots$
  \begin{sousq}
    \item Vérifier que $b = \dfrac{1122}{999}$.
    \item Déterminer le PGCD de 1122 et 999.
    \item Écrire sous forme d'une fraction irréductible le nombre $b$.
    \item Déterminer le PPCM de 1122 et 999.
  \end{sousq}
\end{questions}

\exo{2}{8 points}
Soit la fonction numérique $f$ de la variable réelle $x$ définie par :
\[ f(x) = \frac{x + 3}{-x + 2} \]
On désigne par $(\mathcal{C})$ la courbe représentative de $f$ dans le repère orthonormé $(O,\vi,\vj)$ d'unité graphique 1 cm.
\begin{questions}
  \item Déterminer l'ensemble de définition de $f$.
  \item Déterminer les réels $a$ et $b$ tels que pour tout réel $x$ différent de 2, $f(x) = a + \dfrac{b}{-x + 2}$.
  \item Calculer les limites suivantes : $\lim\limits_{x\to 2^+} f(x)$ ; $\lim\limits_{x\to 2^-} f(x)$ ; $\lim\limits_{x\to-\infty} f(x)$ et $\lim\limits_{x\to+\infty} f(x)$.
  \item Montrer que la dérivée de $f$ a pour expression : $f'(x) = \dfrac{5}{(-x + 2)^2}$.
  \item Déterminer le sens de variation de $f$.
  \item Dresser le tableau de variation de $f$.
  \item Démontrer que les droites d'équations respectives $x = 2$ et $y = -1$ sont respectivement asymptote verticale et asymptote horizontale à la courbe $(\mathcal{C})$.
  \item Tracer dans le repère orthonormal $(O,\vi,\vj)$, la courbe $(\mathcal{C})$ et les droites d'équations respectives $x = 2$ et $y = -1$.
\end{questions}

\exo{3}{4 points}
Un pharmacien observe durant les six premiers mois de l'ouverture de son officine le chiffre d'affaires en millions de francs CFA. Le résultat de l'observation est résumé dans le tableau suivant où $X$ désigne le numéro du mois et $Y$ le chiffre d'affaires correspondant :
\begin{center}\renewcommand{\arraystretch}{1.4}
\begin{tabular}{|c|c|c|c|c|c|c|}
\hline
$X$ & 1 & 2 & 3 & 4 & 5 & 6 \\
\hline
$Y$ & 12 & 13 & 15 & 19 & 21 & 22 \\
\hline
\end{tabular}
\end{center}
\begin{questions}
  \item Calculer les coordonnées $\overline{X}$ et $\overline{Y}$ du point moyen $G$.
  \item Démontrer que l'équation de la droite de régression linéaire de cette série statistique par la méthode de Mayer est $(\mathcal{D}) : y = 2{,}46x + 8{,}36$.
  \item En utilisant la droite de régression linéaire $(\mathcal{D}) : y = 2{,}46x + 8{,}36$, calculer une estimation du chiffre d'affaires de cette pharmacie à la fin de la huitième année.
\end{questions}

\end{document}
